% !TeX spellcheck = ru_RU
% ========== ВСЕ НАСТРОЙКИ ИЗ preambula.tex (КРОМЕ ЛИСТИНГОВ) ==========
\usepackage[T1, T2A]{fontenc}
\usepackage[utf8]{inputenc}
\usepackage[english,russian]{babel}
\usepackage{amssymb,amsfonts,amsmath,mathtext,enumerate,float}
\usepackage{pgfplots}
\usepackage{graphicx}
%\usepackage{tocloft}
\usepackage{caption}
\usepackage{longtable}
\usepackage{tempora}
\usepackage{titlesec}
\usepackage{setspace}
\usepackage{geometry}
\usepackage{indentfirst}
\usepackage{pdfpages}
\usepackage{enumerate,letltxmacro}
\usepackage{threeparttable}
\usepackage{flafter}
\usepackage{enumitem}
\usepackage{multirow}
\usepackage{svg}
\usepackage{float}
\usepackage{hhline}
\usepackage{breqn}
\usepackage{array}
\usepackage{arydshln}
\usepackage{mathtools}
\usepackage[figure,table]{totalcount}
\usepackage{lastpage}
\usepackage{hyphenat}
\usepackage{titling}
\usepackage{newfloat}
\setlist{nosep}
\usepackage{chngcntr} % Для привязки к разделам
\usepackage{listing}
\usepackage[nocompress]{cite}
\usepackage{rotating}
\usepackage{rotfloat}
\usepackage{subcaption}
\usepackage{listings}
\usepackage{hyperref}
\hypersetup{hidelinks}
\usepackage{cleveref}

\captionsetup[longtable]{font=small}

%% Убираем вертикальные отступы между элементами оглавления
%\setlength{\cftbeforechapskip}{1em}   % или 10pt
%\setlength{\cftbeforesecskip}{0pt}    % между разделами внутри главы – нет отступа
%\setlength{\cftbeforesubsecskip}{0pt}
%\setlength{\cftbeforesubsubsecskip}{0pt}

%% Иерархические отступы (восстанавливаем)
%\cftsetindents{chapter}{0pt}{1.5em}
%\cftsetindents{section}{1.5em}{2em}
%\cftsetindents{subsection}{3.5em}{2.5em}
%\cftsetindents{subsubsection}{6em}{3em}

% Лидеры (точки) для всех уровней
%\renewcommand{\cftchapleader}{\cftdotfill{\cftdotsep}}
%\renewcommand{\cftsecleader}{\cftdotfill{\cftdotsep}}
%\renewcommand{\cftsubsecleader}{\cftdotfill{\cftdotsep}}
%\renewcommand{\cftsubsubsecleader}{\cftdotfill{\cftdotsep}}

%\setlength{\cftbeforechapskip}{0pt}        % убрать лишний вертикальный отступ перед главой
%\setlength{\cftbeforesecskip}{0pt}
%\setlength{\cftbeforesubsecskip}{0pt}
%\setlength{\cftbeforesubsubsecskip}{0pt}

% Обнуляем отступы строк и ширину места под номер (левый край строки будет совпадать с левым полем)
%\cftsetindents{chapter}{0pt}{1.5em}
%\cftsetindents{section}{0pt}{2em}
%\cftsetindents{subsection}{0pt}{2.5em}
%\cftsetindents{subsubsection}{0pt}{3em}

% Нумерация листингов в пределах раздела
%\counterwithin{lstlisting}{chapter}

\newcommand{\ssr}[1]{\begin{center}
		\Large\bfseries{#1}
		\addcontentsline{toc}{chapter}{#1}\end{center}}
\newcommand{\ssra}[1]{\begin{center}
		\Large\bfseries{#1}}
	
	\makeatletter
	\renewcommand\LARGE{\@setfontsize\LARGE{22pt}{20}}
	\renewcommand\Large{\@setfontsize\Large{20pt}{20}}
	\renewcommand\large{\@setfontsize\large{16pt}{20}}
	\makeatother
	\RequirePackage{titlesec}
	\titleformat{\chapter}[block]{\hspace{\parindent}\large\bfseries}{\thechapter}{0.5em}{\large\bfseries\raggedright}
	\titleformat{name=\chapter,numberless}[block]{\hspace{\parindent}}{}{0pt}{\large\bfseries\centering}
	\titleformat{\section}[block]{\hspace{\parindent}\large\bfseries}{\thesection}{0.5em}{\large\bfseries\raggedright}
	\titleformat{\subsection}[block]{\hspace{\parindent}\large\bfseries}{\thesubsection}{0.5em}{\large\bfseries\raggedright}
	\titleformat{\subsubsection}[block]{\hspace{\parindent}\large\bfseries}{\thesubsection}{0.5em}{\large\bfseries\raggedright}
	\titlespacing{\chapter}{12.5mm}{-22pt}{10pt}
	\titlespacing{\section}{12.5mm}{10pt}{10pt}
	\titlespacing{\subsection}{12.5mm}{10pt}{10pt}
	\titlespacing{\subsubsection}{12.5mm}{10pt}{10pt}
	
	\makeatletter
	\renewcommand{\@biblabel}[1]{#1.}
	\makeatother
	
	\titleformat{\chapter}[hang]{\LARGE\bfseries}{\hspace{1.25cm}\thechapter}{1ex}{\LARGE\bfseries}
	\titleformat{\section}[hang]{\Large\bfseries}{\hspace{1.25cm}\thesection}{1ex}{\Large\bfseries}
	\titleformat{name=\section,numberless}[hang]{\Large\bfseries}{\hspace{1.25cm}}{0pt}{\Large\bfseries}
	\titleformat{\subsection}[hang]{\large\bfseries}{\hspace{1.25cm}\thesubsection}{1ex}{\large\bfseries}
	\titlespacing{\chapter}{0pt}{-\baselineskip}{\baselineskip}
	\titlespacing*{\section}{0pt}{\baselineskip}{\baselineskip}
	\titlespacing*{\subsection}{0pt}{\baselineskip}{\baselineskip}
	
	\geometry{left=30mm}
	\geometry{right=10mm}
	\geometry{top=20mm}
	\geometry{bottom=20mm}
	
	\onehalfspacing
	
	\renewcommand{\theenumi}{\arabic{enumi}}
	\renewcommand{\labelenumi}{\arabic{enumi}\text{)}}
	\renewcommand{\theenumii}{.\arabic{enumii}}
	\renewcommand{\labelenumii}{\asbuk{enumii}\text{)}}
	\renewcommand{\theenumiii}{.\arabic{enumiii}}
	\renewcommand{\labelenumiii}{\arabic{enumi}.\arabic{enumii}.\arabic{enumiii}.}
	
	%\renewcommand{\cftchapleader}{\cftdotfill{\cftdotsep}}
	
	
	\addto\captionsrussian{\renewcommand{\figurename}{Рисунок}}
	\DeclareCaptionLabelSeparator{dash}{~---~}
	\captionsetup{labelsep=dash}
	\captionsetup[figure]{justification=centering,labelsep=dash}
	\captionsetup[table]{labelsep=dash,justification=raggedright,singlelinecheck=off}
	\captionsetup[lstlisting]{justification=raggedright,singlelinecheck=off}
	
	\graphicspath{{images/}}
	
	\newcommand{\mytexttilde}{\raisebox{0.5ex}{\texttildelow}}
	\newcommand{\floor}[1]{\lfloor #1 \rfloor}
	
	\pgfplotsset{width=0.85\linewidth, height=0.5\columnwidth}
	\linespread{1.3}
	\parindent=1.25cm
	
	\def\labelitemi{---}
	\setlist[itemize]{leftmargin=1.25cm, itemindent=0.65cm}
	\setlist[enumerate]{leftmargin=1.25cm, itemindent=0.55cm}
	\newcommand{\specialcell}[2][c]{%
		\begin{tabular}[#1]{@{}c@{}}#2\end{tabular}}
	\frenchspacing
	
	% ========== НАСТРОЙКИ ДЛЯ ЛИСТИНГОВ И ПСЕВДОКОДА (ИЗ settings.tex) ==========
	\usepackage{cmap}
	\usepackage{csquotes}
	\usepackage{url}
	\usepackage{tikz}
	\usepackage{booktabs}
	\usepackage{ulem}
	\usepackage{verbatim}
	\usepackage{xcolor}
	
	% Дополнительные настройки из settings.tex (не дублируем уже загруженные)
	\newenvironment{signstabular}[1][1]{
		\renewcommand*{\arraystretch}{#1}
		\tabular
	}{
		\endtabular
	}
	
	%\setenumerate[0]{label=\arabic*.)}
	\renewcommand{\labelitemi}{---}
	
	\NewDocumentCommand{\ulinetext}{O{3cm} O{c} m m}{
		$\underset{\text{#3}}{\uline{\makebox[#1][#2]{#4}}}$
	}
	
	%\makeatletter
	%\renewcommand*{\l@chapter}[2]{
		%	\ifnum \c@tocdepth>\m@ne
		%	\addpenalty{-\@highpenalty}
		%	\vskip 1em \@plus.2em
		%	\@dottedtocline{0}{0pt}{1.5em}{\bfseries #1}{\bfseries #2}
		%	\fi
		%}
	%\renewcommand*{\l@subsubsection}[2]{
		%	\ifnum \c@tocdepth>2
		%	\@dottedtocline{3}{6.9em}{3.8em}{#1}{#2}
		%	\fi
		%}
	%\makeatother
	
	\makeatletter
	\renewcommand*{\l@chapter}[2]{
		\ifnum \c@tocdepth>\m@ne
		\addpenalty{-\@highpenalty}
		\vskip 1em \@plus.2em
		\@dottedtocline{0}{0pt}{1.5em}{\bfseries #1}{\bfseries #2}
		\fi
	}
	\renewcommand*{\l@section}[2]{\@dottedtocline{1}{1.5em}{2em}{#1}{#2}}
	\renewcommand*{\l@subsection}[2]{\@dottedtocline{2}{3.5em}{2.5em}{#1}{#2}}
	\renewcommand*{\l@subsubsection}[2]{\@dottedtocline{3}{6em}{3em}{#1}{#2}}
	\makeatother
	
	\setcounter{tocdepth}{3}
	\setcounter{secnumdepth}{3}
	
	% Настройки отображения кода (listings)
	\lstset{
		basicstyle=\footnotesize\ttfamily,
		language=SQL,
		numbers=left,
		numbersep=5pt,
		stepnumber=1,
		xleftmargin=17pt,
		frame=single,
		tabsize=4,
		captionpos=t,
		breaklines=true,
		breakatwhitespace=true,
		escapeinside={\#*}{*)},
		inputencoding=utf8x,
		backgroundcolor=\color{white},
		numberstyle=\scriptsize,
		keywordstyle=\color{black},
		stringstyle=\color{black},
		commentstyle=\color{black},
		morekeywords={TO,WITH,SUPERUSER,PRIVILEGES,TABLES,SCHEMA,USER,PASSWORD,public,REFERENCES,UUID,IF}
	}
	\lstset{
		literate=
		{~}{{\textasciitilde}}1
		{а}{{\selectfont\char224}}1
		{б}{{\selectfont\char225}}1
		{в}{{\selectfont\char226}}1
		{г}{{\selectfont\char227}}1
		{д}{{\selectfont\char228}}1
		{е}{{\selectfont\char229}}1
		{ё}{{\"e}}1
		{ж}{{\selectfont\char230}}1
		{з}{{\selectfont\char231}}1
		{и}{{\selectfont\char232}}1
		{й}{{\selectfont\char233}}1
		{к}{{\selectfont\char234}}1
		{л}{{\selectfont\char235}}1
		{м}{{\selectfont\char236}}1
		{н}{{\selectfont\char237}}1
		{о}{{\selectfont\char238}}1
		{п}{{\selectfont\char239}}1
		{р}{{\selectfont\char240}}1
		{с}{{\selectfont\char241}}1
		{т}{{\selectfont\char242}}1
		{у}{{\selectfont\char243}}1
		{ф}{{\selectfont\char244}}1
		{х}{{\selectfont\char245}}1
		{ц}{{\selectfont\char246}}1
		{ч}{{\selectfont\char247}}1
		{ш}{{\selectfont\char248}}1
		{щ}{{\selectfont\char249}}1
		{ъ}{{\selectfont\char250}}1
		{ы}{{\selectfont\char251}}1
		{ь}{{\selectfont\char252}}1
		{э}{{\selectfont\char253}}1
		{ю}{{\selectfont\char254}}1
		{я}{{\selectfont\char255}}1
		{А}{{\selectfont\char192}}1
		{Б}{{\selectfont\char193}}1
		{В}{{\selectfont\char194}}1
		{Г}{{\selectfont\char195}}1
		{Д}{{\selectfont\char196}}1
		{Е}{{\selectfont\char197}}1
		{Ё}{{\"E}}1
		{Ж}{{\selectfont\char198}}1
		{З}{{\selectfont\char199}}1
		{И}{{\selectfont\char200}}1
		{Й}{{\selectfont\char201}}1
		{К}{{\selectfont\char202}}1
		{Л}{{\selectfont\char203}}1
		{М}{{\selectfont\char204}}1
		{Н}{{\selectfont\char205}}1
		{О}{{\selectfont\char206}}1
		{П}{{\selectfont\char207}}1
		{Р}{{\selectfont\char208}}1
		{С}{{\selectfont\char209}}1
		{Т}{{\selectfont\char210}}1
		{У}{{\selectfont\char211}}1
		{Ф}{{\selectfont\char212}}1
		{Х}{{\selectfont\char213}}1
		{Ц}{{\selectfont\char214}}1
		{Ч}{{\selectfont\char215}}1
		{Ш}{{\selectfont\char216}}1
		{Щ}{{\selectfont\char217}}1
		{Ъ}{{\selectfont\char218}}1
		{Ы}{{\selectfont\char219}}1
		{Ь}{{\selectfont\char220}}1
		{Э}{{\selectfont\char221}}1
		{Ю}{{\selectfont\char222}}1
		{Я}{{\selectfont\char223}}1
	}
	
	% Настройка листингов с tcolorbox и minted
	\RequirePackage[listings]{tcolorbox}
	\tcbuselibrary{listings,breakable,minted,skins}
	\renewcommand{\theFancyVerbLine}{\scriptsize\arabic{FancyVerbLine}}
	
	\newtcbinputlisting[auto counter,number within=chapter]{\mylisting}[4][]{
		enhanced,
		before skip=10pt,
		arc=0mm,
		boxrule=1pt,
		listing only,
		listing file={chapters/#1},
		listing engine=minted,
		minted language=python,
		minted options = {
			linenos,
			breaklines=true,
			breakbefore=.,
			fontsize=\small\ttfamily,
			numbersep=.1mm,
			firstnumber=1,
			escapeinside=~~,
			style = algol,
			#3
		},
		colback=white,
		colframe=black,
		breakable,
		title={\parbox{\linewidth}{Листинг~\thetcbcounter~--~#2}},
		attach boxed title to top left,
		boxed title style={enhanced jigsaw, opacityfill=1, sharp corners, boxrule=0pt, left=-3pt, right=0pt},
		coltitle=black,
		colbacktitle=white,
		enlarge top at break by=5mm,
		overlay first={
			\draw[line width=.5pt] (frame.south west)--(frame.south east);
		},
		overlay middle={
			\draw[line width=.5pt] (frame.south west)--(frame.south east);
			\draw[line width=.5pt] (frame.north west)--(frame.north east);
			\node[anchor=south west, inner xsep=0pt, xshift=0pt] at (frame.north west) {\parbox{\linewidth}{Продолжение листинга \thetcbcounter~--~#2}};
		},
		overlay last={
			\draw[line width=.5pt] (frame.north west)--(frame.north east);
			\node[anchor=south west, inner xsep=0pt, xshift=0pt] at (frame.north west) {\parbox{\linewidth}{Продолжение листинга \thetcbcounter~--~#2}};
		},
		label=lst:#1,
		#4
	}
	
	% Для приложений
	\newcounter{mylistingappcounter}
	\newcommand{\appendixletter}{A}
	\newtcbinputlisting[use counter=mylistingappcounter, list inside=lol, list type=lol]{\mylistingapp}[4][]{
		enhanced,
		before skip=10pt,
		arc=0mm,
		boxrule=1pt,
		listing only,
		listing file={inc/lst/#1},
		listing engine=minted,
		minted language=C,
		minted options = {
			linenos,
			breaklines=true,
			breakbefore=.,
			fontsize=\small\ttfamily,
			numbersep=.1mm,
			firstnumber=1,
			escapeinside=~~,
			#3
		},
		colback=white,
		colframe=black,
		breakable,
		title={Листинг~\appendixletter.\thetcbcounter~--~#2},
		attach boxed title to top left,
		boxed title style={enhanced jigsaw, opacityfill=1, sharp corners, boxrule=0pt, left=-3pt, right=0pt},
		coltitle=black,
		colbacktitle=white,
		enlarge top at break by=5mm,
		overlay first={
			\draw[line width=.5pt] (frame.south west)--(frame.south east);
		},
		overlay middle={
			\draw[line width=.5pt] (frame.south west)--(frame.south east);
			\draw[line width=.5pt] (frame.north west)--(frame.north east);
			\node[anchor=south west, inner xsep=0pt, xshift=0pt] at (frame.north west) {Продолжение листинга \appendixletter.\thetcbcounter};
		},
		overlay last={
			\draw[line width=.5pt] (frame.north west)--(frame.north east);
			\node[anchor=south west, inner xsep=0pt, xshift=0pt] at (frame.north west) {Окончание листинга \appendixletter.\thetcbcounter};
		},
		label=lst:#1,
		#4
	}
	
	% Окружение для псевдокода алгоритмов
	\usepackage[ruled,vlined,linesnumbered,noline]{algorithm2e}
	\SetAlgorithmName{Листинг}{листинг}{Список листингов}
	\makeatletter
	\renewcommand{\algocf@caption@ruled}{}
	\renewcommand{\algocf@caption@plain}{}
	\makeatother
	\SetKwInput{Input}{Вход}
	\SetKwInput{Output}{Выход}
	\SetKw{Return}{вернуть}
	\SetKwBlock{Begin}{начало}{конец}
	\SetKwIF{If}{ElseIf}{Else}{если}{тогда}{иначе если}{иначе}{конец если}
	\SetKwFor{For}{для}{выполнять}{конец цикла}
	\SetKwFor{ForEach}{для каждого}{выполнять}{конец цикла}
	\SetKwFor{While}{пока}{выполнять}{конец цикла}
	\SetNlSty{textbf}{}{}
	\SetAlgoNlRelativeSize{-1}
	\SetNlSkip{0.8em}
	\SetCommentSty{textit}
	\SetKwComment{Comment}{$\triangleright$~}{}
	\RestyleAlgo{boxed}
	\SetAlgoSkip{smallskip}
	\newcounter{myalgocounter}
	\counterwithin{myalgocounter}{chapter}
	\renewcommand{\themyalgocounter}{\thechapter.\arabic{myalgocounter}}
	\newcommand{\algocaption}[2]{%
		\refstepcounter{myalgocounter}%
		\label{#1}%
		\noindent Листинг~\themyalgocounter~-- #2\par%
		\vspace{2pt}%
	}
	\usepackage{etoolbox}
	\AfterEndEnvironment{algorithm}{\vspace{\baselineskip}}
	\usepackage{totcount}
	\regtotcounter{figure}
	\regtotcounter{table}
	\newtotcounter{citnum}
	%\AtEveryBibitem{\stepcounter{citnum}}
	
	%\usepackage{appendix}                     % для удобной работы с приложениями
	%\makeatletter
	%\newcommand{\appendixsetup}{%
		%	% Стандартные плавающие объекты
		%	\renewcommand{\thetable}{\thechapter.\arabic{table}}%
		%	\renewcommand{\thefigure}{\thechapter.\arabic{figure}}%
		%	\renewcommand{\thelstlisting}{\thechapter.\arabic{lstlisting}}%
		%	% Алгоритмы (пакет algorithm2e)
		%	\renewcommand{\thealgocf}{\thechapter.\arabic{algocf}}%
		%	% Листинги через tcolorbox (счётчик tcbcounter)
		%	\renewcommand{\thetcbcounter}{\thechapter.\arabic{tcbcounter}}%
		%	% Ваш счётчик для псевдокода
		%	\renewcommand{\themyalgocounter}{\thechapter.\arabic{myalgocounter}}%
		%	% Сброс счётчиков при начале каждой главы приложения
		%	\counterwithin{table}{chapter}%
		%	\counterwithin{figure}{chapter}%
		%	\counterwithin{lstlisting}{chapter}%
		%	\counterwithin{algocf}{chapter}%
		%	\counterwithin{tcbcounter}{chapter}%
		%	\counterwithin{myalgocounter}{chapter}%
		%}
	%% Автоматически вызываем настройку после команды \appendix
	%\let\originalappendix\appendix
	%\renewcommand{\appendix}{%
		%	\originalappendix
		%	\appendixsetup
		%}
	%\makeatother
	
	\newcommand{\startappendix}[2]{%
		\ssr{#1 #2}%
		% Буква приложения теперь задаётся вторым аргументом
		\renewcommand{\thetable}{#2.\arabic{table}}%
		\renewcommand{\thefigure}{#2.\arabic{figure}}%
		\renewcommand{\thelstlisting}{#2.\arabic{lstlisting}}%
		\renewcommand{\thealgocf}{#2.\arabic{algocf}}%
		\ifcsname thetcbcounter\endcsname
		\renewcommand{\thetcbcounter}{#2.\arabic{tcbcounter}}%
		\fi
		\ifcsname themyalgocounter\endcsname
		\renewcommand{\themyalgocounter}{#2.\arabic{myalgocounter}}%
		\fi
		% Сброс счётчиков
		\setcounter{table}{0}%
		\setcounter{figure}{0}%
		\setcounter{lstlisting}{0}%
		\setcounter{algocf}{0}%
		\ifcsname c@tcbcounter\endcsname \setcounter{tcbcounter}{0}\fi
		\ifcsname c@myalgocounter\endcsname \setcounter{myalgocounter}{0}\fi
	}